% !PRINTED 381
\weberpage{381}
%<<CONT
of degree $4$, $F[x, \Theta(x-\xi)]$, be a complete square. This gives two equations between $\xi$ and $\Theta$, from which, by elimination of $\Theta$, one obtains the double-tangent equation. Since to each $\xi$ there belongs only one value of $\Theta$ (so long as no two double tangents intersect on the line $L$), $\Theta$ can be expressed rationally in terms of $\xi$, and indeed in a form
\weq{2}{\Theta = \varphi\,(\xi),}
which holds for every pair $\xi$, $\Theta$ that belong together.

The roots of the double-tangent equation arrange themselves, just as the corresponding double tangents do, into complexes, seven-systems, and so on. We denote these roots, just like the double tangents in \S.\,99, by the symbol $[ik]$, in which $i$, $k$ run through the pairs of the figures 1 to 8 and $[ik] = [ki]$.

If we regard any two of the double tangents, $p$, $q$, as known, then we can derive from them by a rational path the equation $u = 0$ of a conic which passes through the four points of contact of these double tangents, and we can therefore set up the equation \S.\,96, (2)
\[
f = p\,q\,v - u^2
\]
in rational form. Then, by \S.\,96, (4), there are among the pencil of conics
\[
v + 2\,\lambda\,u + \lambda^2\,p\,q = 0
\]
five which split into a pair of lines, and if we form the product
\[
\prod_{\lambda} (v + 2\,\lambda\,u + \lambda^2\,p\,q) = \Phi
\]
over the five roots of the equation of degree $5$ on which $\lambda$ depends [\S.\,96, (6)], then this product is likewise expressible rationally in terms of the coefficients of $f$ and of $p$, $q$. But this product is a form of degree $10$, which splits into ten linear factors, which together with $p\,q$ form a \gesperrt{Steiner} complex.

The coefficients in $\Phi$ can now also be expressed rationally in terms of the two roots $\xi_1$, $\xi_2$ of the double-tangent equation corresponding to $p$, $q$, and if we then seek the abscissae of the points of intersection of the line $L$ with $\Phi = 0$, we obtain an equation of degree $10$ for $\xi$
\[
X(\xi_1, \xi_2, \xi) = 0,
\]
%CONT>>
